\section{Basic operations and linear convolution (2.4--2.15)}

\Prob{2.4\quad Input--output relations for the four structures of Figure P2.1}

\textbf{(a)} Let $v[n]$ be the signal at the node immediately to the right of the input adder.
The left half is a transposed feedback section: the upper delay feeds the input adder, and the
input of that delay is the sum of $-d_1v[n]$ and the output of the lower delay, whose input is
$-d_2v[n]$. With the two state variables $s_1,s_2$,
\[
s_1[n]=-d_1v[n-1]+s_2[n-1],\qquad s_2[n]=-d_2v[n-1],\qquad v[n]=x[n]+s_1[n],
\]
so that
\[
v[n]=x[n]-d_1v[n-1]-d_2v[n-2].
\]
The right half taps $v[n]$ through $p_0,p_1,p_2$ into a transposed feed-forward chain, giving
$y[n]=p_0v[n]+p_1v[n-1]+p_2v[n-2]$. Eliminating $v$ (in the $z$-domain,
$V(z)=X(z)/(1+d_1z^{-1}+d_2z^{-2})$ and $Y(z)=(p_0+p_1z^{-1}+p_2z^{-2})V(z)$):
\[
\boxed{\;y[n]+d_1\,y[n-1]+d_2\,y[n-2]=p_0\,x[n]+p_1\,x[n-1]+p_2\,x[n-2].\;}
\]

\textbf{(b)} Here the feed-forward chain acts directly on $x[n]$ and the feedback chain
directly on $y[n]$:
\[
y[n]=\bigl(p_0x[n]+p_1x[n-1]+p_2x[n-2]\bigr)+\bigl(-d_1y[n-1]-d_2y[n-2]\bigr),
\]
that is,
\[
\boxed{\;y[n]+d_1y[n-1]+d_2y[n-2]=p_0x[n]+p_1x[n-1]+p_2x[n-2].\;}
\]
\textbf{Structures (a) and (b) realize the same difference equation}; they are two
realizations of
$H(z)=\dfrac{p_0+p_1z^{-1}+p_2z^{-2}}{1+d_1z^{-1}+d_2z^{-2}}$.

\textbf{(c)} A gain $h[0]$ followed by three identical second-order feed-forward sections.
Section $i$ adds to its own input that input delayed once (weight $\beta_{1i}$) and delayed
twice (weight $\beta_{2i}$), so it has transfer function $1+\beta_{1i}z^{-1}+\beta_{2i}z^{-2}$.
Hence
\[
\boxed{\;Y(z)=h[0]\prod_{i=1}^{3}\bigl(1+\beta_{1i}z^{-1}+\beta_{2i}z^{-2}\bigr)X(z),\;}
\]
equivalently, in the time domain,
\[
y[n]=h[0]\;x[n]\circledast g_1[n]\circledast g_2[n]\circledast g_3[n],\qquad
g_i[n]=\delta[n]+\beta_{1i}\delta[n-1]+\beta_{2i}\delta[n-2].
\]
This is a \emph{cascade} realization of a length-7 FIR filter.

\textbf{(d)} The upper delay line produces $x[n],x[n-1],x[n-2],x[n-3]$; the sample $x[n-3]$ is
then fed into the lower delay line, which runs right-to-left and produces
$x[n-4],x[n-5],x[n-6]$. The three adders pair taps symmetrically about $x[n-3]$, and only
$x[n-3]$ is unpaired, so
\[
\boxed{\;
\begin{aligned}
y[n]=\;&h[0]\bigl(x[n]+x[n-6]\bigr)+h[1]\bigl(x[n-1]+x[n-5]\bigr)\\
      &+h[2]\bigl(x[n-2]+x[n-4]\bigr)+h[3]\,x[n-3].
\end{aligned}\;}
\]
This is the 4-multiplier realization of a length-7 \emph{symmetric} (Type-I linear-phase)
FIR filter: 6 delays and 4 multipliers instead of 7.

\Prob{2.5\quad The convolution of a length-$M$ with a length-$N$ sequence has length $M+N-1$}

Take $x[n]$ nonzero only for $0\le n\le M-1$ and $h[n]$ only for $0\le n\le N-1$ (shifting
either sequence only shifts the result, it does not change its length). In
\[
y[n]=\sum_{k}x[k]h[n-k]
\]
a summand can be nonzero only if some $k$ satisfies simultaneously
\[
0\le k\le M-1 \qquad\text{and}\qquad 0\le n-k\le N-1 .
\]
Adding the two smallest admissible values gives $n\ge 0$; adding the two largest gives
$n\le (M-1)+(N-1)=M+N-2$. Outside that window the two index intervals do not overlap and
$y[n]=0$. The end samples
\[
y[0]=x[0]h[0]\ne 0,\qquad y[M+N-2]=x[M-1]h[N-1]\ne 0
\]
are nonzero, because the first and last samples of a length-$M$ (length-$N$) sequence are
nonzero by definition. Hence the support is exactly $0\le n\le M+N-2$ and
\[
\boxed{\;\text{length}(y)=(M+N-2)-0+1=M+N-1.\;}
\]

\Prob{2.6\quad Length of $x[n]\circledast y[n]\circledast w[n]$}

Apply Problem 2.5 twice. $x\circledast y$ has length $N+M-1$; convolving that with $w$
(length $L$) gives
\[
\boxed{\;(N+M-1)+L-1=N+M+L-2 .\;}
\]
Since all three sequences start at $n=0$, the result is defined for $0\le n\le N+M+L-3$.

\Prob{2.7\quad Show that $x_1[n]\circledast x_2[n]=x_3[n]\circledast x_4[n]$}

The clean route is the polynomial ($z$-domain) one. With
$\{x_1[n]\}=\{1,0,1,0,\dots,1\}$ (15 samples: ones at the even indices $0,2,\dots,14$) and
$\{x_2[n]\}=\{1,1\}$,
\[
X_1(z)=\sum_{m=0}^{7}z^{-2m}=\frac{1-z^{-16}}{1-z^{-2}},\qquad X_2(z)=1+z^{-1},
\]
\[
X_1(z)X_2(z)=\frac{(1-z^{-16})(1+z^{-1})}{(1-z^{-1})(1+z^{-1})}
            =\frac{1-z^{-16}}{1-z^{-1}}=\sum_{n=0}^{15}z^{-n}.
\]
Hence
\[
\boxed{\;x_1[n]\circledast x_2[n]=\{\underbrace{1,1,\dots,1}_{16}\},\qquad 0\le n\le 15.\;}
\]
For the second pair, $\{x_3[n]\}=\{1,1,0,0,0,0,0,0,1,1\}$ gives
\[
X_3(z)=1+z^{-1}+z^{-8}+z^{-9}=(1+z^{-1})(1+z^{-8}).
\]
Using the standard factorization
\[
\frac{1-z^{-16}}{1-z^{-1}}=(1+z^{-1})(1+z^{-2})(1+z^{-4})(1+z^{-8}),
\]
the identity $X_3X_4=X_1X_2$ forces
\[
X_4(z)=(1+z^{-2})(1+z^{-4})=1+z^{-2}+z^{-4}+z^{-6}
\;\Longleftrightarrow\;
\boxed{\;\{x_4[n]\}=\{1,\,0,\,1,\,0,\,1,\,0,\,1\}.\;}
\]
With this $x_4$ the required equality holds, both sides being 16 consecutive ones.

\smallskip
\Mis The scanned page lists $\{x_4[n]\}=\{1,1,0,0,0,1,1\}$, i.e.\
$X_4(z)=(1+z^{-1})(1+z^{-5})$. With \emph{that} $x_4$,
\[
x_3[n]\circledast x_4[n]=\{1,2,1,0,0,1,2,1,1,2,1,0,0,1,2,1\}\;\ne\;x_1[n]\circledast x_2[n],
\]
so the stated identity is false as printed. The intended sequence is
$\{1,0,1,0,1,0,1\}$ (checked numerically in \texttt{hw1\_p207.m}).

\Prob{2.8\quad Linear convolution of each sequence with itself}

\textbf{(a)} $x_1[n]=\{1,-1,1\}$, $-1\le n\le 1$. Here $X_1(z)=z-1+z^{-1}$, so
$X_1^2(z)=z^{2}-2z+3-2z^{-1}+z^{-2}$:
\[
\boxed{\;x_1\circledast x_1=\{1,\,-2,\,3,\,-2,\,1\},\qquad -2\le n\le 2.\;}
\]

\textbf{(b)} $x_2[n]=\{1,-1,0,1,-1\}$ starting at $n=0$. Note
$X_2(z)=(1-z^{-1})(1+z^{-3})$, hence
$X_2^2(z)=(1-2z^{-1}+z^{-2})(1+2z^{-3}+z^{-6})$:
\[
\boxed{\;x_2\circledast x_2=\{1,\,-2,\,1,\,2,\,-4,\,2,\,1,\,-2,\,1\},\qquad 0\le n\le 8.\;}
\]
\Mis The page states the range $0\le n\le 5$ but lists only five samples; it should read
$0\le n\le 4$. (Appending a trailing zero to make six samples changes none of the values
above, only the nominal length.)

\textbf{(c)} $x_3[n]=\{-1,2,0,-2,1\}$, $-3\le n\le 1$:
\[
\boxed{\;x_3\circledast x_3=\{1,\,-4,\,4,\,4,\,-10,\,4,\,4,\,-4,\,1\},\qquad -6\le n\le 2.\;}
\]
Consistency check: $x_3[n]$ is antisymmetric about $n=-1$, so $x_3\circledast x_3$ must be
symmetric about $n=-2$ --- and it is.

\Prob{2.9\quad Convolutions of the Problem-2.3 sequences}

The starting index of a linear convolution is the sum of the starting indices, and the length
is the sum of the lengths minus one (Problem 2.5). All three sequences have length 7, so each
result has length 13.

\textbf{(a)} $u[n]=x[n]\circledast y[n]$, starting index $-1+(-3)=-4$:
\[
\boxed{\;\{u[n]\}=\{-15,\,17,\,14,\,-14,\,1,\,-41,\,49,\,5,\,-23,\,37,\,-20,\,14,\,-24\},\quad -4\le n\le 8.\;}
\]

\textbf{(b)} $v[n]=x[n]\circledast w[n]$, starting index $-1+3=2$:
\[
\boxed{\;\{v[n]\}=\{-3,\,13,\,-22,\,5,\,26,\,-40,\,42,\,-22,\,-25,\,31,\,-8,\,21,\,-18\},\quad 2\le n\le 14.\;}
\]

\textbf{(c)} $g[n]=w[n]\circledast y[n]$, starting index $3+(-3)=0$:
\[
\boxed{\;\{g[n]\}=\{5,\,-14,\,17,\,1,\,-2,\,32,\,-63,\,29,\,-41,\,-14,\,54,\,-52,\,48\},\quad 0\le n\le 12.\;}
\]
Check: $\sum_n u[n]=\bigl(\sum_n x[n]\bigr)\bigl(\sum_n y[n]\bigr)=(0)(9)=0$, and the thirteen
values of $u[n]$ do sum to zero. Likewise $\sum_n v[n]=(0)(0)=0$ and
$\sum_n g[n]=(0)(9)=0$.

\Prob{2.10\quad $v[n]=x_1[n-N_1]\circledast x_2[n-N_2]$ in terms of $y[n]=x_1[n]\circledast x_2[n]$}

\[
v[n]=\sum_{k}x_1[k-N_1]\,x_2[n-k-N_2]
     \;\overset{m=k-N_1}{=}\;\sum_{m}x_1[m]\,x_2\bigl((n-N_1-N_2)-m\bigr),
\]
which is $y[\cdot]$ evaluated at $n-N_1-N_2$:
\[
\boxed{\;v[n]=y[n-N_1-N_2].\;}
\]
Equivalently, in the $z$-domain each shift contributes a factor $z^{-N_i}$, so
$V(z)=z^{-(N_1+N_2)}Y(z)$: convolution turns individual delays into one accumulated delay.

\Prob{2.11\quad Three-fold version of Problem 2.10}

The same substitution applied to each factor (or multiplying the three delay factors
$z^{-N_1}z^{-N_2}z^{-N_3}$) gives
\[
\boxed{\;h[n]=g[n-N_1-N_2-N_3].\;}
\]

\Prob{2.12\quad Lengths and index ranges}

The given sequences have
\[
\begin{aligned}
h[n]:&\;-M\le n\le N &&\Rightarrow\; \text{length } N+M+1,\\
g[n]:&\;\phantom{-}K\le n\le N &&\Rightarrow\; \text{length } N-K+1,\\
w[n]:&\;-L\le n\le -R &&\Rightarrow\; \text{length } L-R+1 .
\end{aligned}
\]
For a convolution the index range runs from the sum of the lower limits to the sum of the
upper limits, and the length is the sum of the lengths minus one:

\begin{center}
\renewcommand{\arraystretch}{1.3}
\begin{tabular}{@{}llll@{}}
\toprule
& sequence & length & range of $n$\\
\midrule
(a) & $y_1=h\circledast h$ & $2(N+M)+1$ & $-2M\le n\le 2N$\\
(b) & $y_2=g\circledast g$ & $2(N-K)+1$ & $2K\le n\le 2N$\\
(c) & $y_3=w\circledast w$ & $2(L-R)+1$ & $-2L\le n\le -2R$\\
(d) & $y_4=h\circledast g$ & $2N+M-K+1$ & $K-M\le n\le 2N$\\
(e) & $y_5=h\circledast w$ & $N+M+L-R+1$ & $-M-L\le n\le N-R$\\
\bottomrule
\end{tabular}
\end{center}
The hypotheses $K<N$ and $L>R$ guarantee that $g$ and $w$ have length at least 2.

\Prob{2.13\quad $y[-1]$ without computing the convolution sum}

$\{x[n]\}=\{2,-3,4,1\}$ for $-1\le n\le 2$ and $\{h[n]\}=\{-3,5,-6,4\}$ for $-2\le n\le 1$.
In
\[
y[-1]=\sum_k x[k]\,h[-1-k]
\]
the summand is nonzero only when $-1\le k\le 2$ \emph{and} $-2\le -1-k\le 1$, i.e.\
$-2\le k\le 1$. The intersection is $k\in\{-1,0,1\}$, so only three products are needed:
\[
y[-1]=x[-1]h[0]+x[0]h[-1]+x[1]h[-2]=(2)(-6)+(-3)(5)+(4)(-3)=-12-15-12 .
\]
\[
\boxed{\;y[-1]=-39.\;}
\]
(The full convolution is $\{-6,\,19,\,-39,\,43,\,-31,\,10,\,4\}$ for $-3\le n\le 3$; its third
sample is indeed $-39$.)

\Prob{2.14\quad Largest positive sample of $y[n]=x[n]\circledast h[n]$ for two ramps}

Here $x[n]=n$ for $0\le n\le N-1$ and $h[n]=n$ for $0\le n\le N-3$, so
\[
y[n]=\sum_{k=a(n)}^{b(n)}k\,(n-k),\qquad a(n)=\max(0,\,n-N+3),\quad b(n)=\min(N-1,\,n),
\]
defined for $0\le n\le 2N-4$ (length $2N-3$).

\textbf{Closed form.} Using $\sum_{k=a}^{b}k=\tfrac{b(b+1)-(a-1)a}{2}$ and
$\sum_{k=a}^{b}k^{2}=\tfrac{b(b+1)(2b+1)-(a-1)a(2a-1)}{6}$,
\[
y[n]=n\cdot\frac{b(b+1)-(a-1)a}{2}-\frac{b(b+1)(2b+1)-(a-1)a(2a-1)}{6}.
\]
In particular, in the initial region $0\le n\le N-3$ (where $a=0,\;b=n$),
\[
y[n]=n\cdot\frac{n(n+1)}{2}-\frac{n(n+1)(2n+1)}{6}=\frac{n^{3}-n}{6},
\]
which is strictly increasing --- so the maximum lies in the trailing region.

\textbf{Location of the maximum, without evaluating the sum.} For $n\ge N-2$ we have
$b=N-1$ and $a=n-N+3$; the forward difference telescopes to
\[
\Delta y[n]=y[n+1]-y[n]=\sum_{k=n-N+4}^{N-1}k\;-\;(n-N+3)(N-3).
\]
The first term decreases in $n$ and the second increases, so $\Delta y[n]$ changes sign
exactly once: $y[n]$ is unimodal. Setting $\Delta y=0$ with $m\triangleq n-N+3$,
\[
\frac{(N-1)N-m(m+1)}{2}=m(N-3)
\;\Longrightarrow\;
m^{2}+(2N-5)\,m-N(N-1)=0,
\]
\[
\boxed{\;m^\star=\frac{5-2N+\sqrt{8N^{2}-24N+25}}{2},\qquad
n_{\max}=\bigl\lceil\, m^\star\,\bigr\rceil+N-3 .\;}
\]
The corresponding value is $y[n_{\max}]$ from the closed form above. For large $N$,
$m^\star\to(\sqrt2-1)N$ and therefore $n_{\max}\approx\sqrt{2}\,N$: the peak sits at about
$71\%$ of the total length $2N-4$, not at the midpoint --- a consequence of the ramps
weighting the late samples most heavily.

\textbf{Examples} (confirmed by direct convolution in \texttt{hw1\_p214.m}):
\begin{center}
\renewcommand{\arraystretch}{1.15}
\begin{tabular}{@{}rrrr@{}}
\toprule
$N$ & $m^\star$ & $n_{\max}$ & $y[n_{\max}]$\\
\midrule
$5$ & $2.6235$ & $5$ & $10$\\
$6$ & $3.0000$ & $6$ & $22$ \ \ (tie with $n=7$)\\
$8$ & $3.7871$ & $9$ & $72$\\
$10$ & $4.5934$ & $12$ & $165$\\
$15$ & $6.6377$ & $19$ & $672$\\
$20$ & $8.6963$ & $26$ & $1738$\\
\bottomrule
\end{tabular}
\end{center}
When $m^\star$ happens to be an integer (as for $N=6$) the difference $\Delta y$ vanishes
exactly and the maximum is attained at two adjacent indices.

\Prob{2.15\quad Recovering $x[n]$ from $y[n]=h[n]\circledast x[n]$ (deconvolution)}

Because both sequences are causal, the convolution sum can be solved recursively for
$x[n]$ (polynomial long division of $Y(z)$ by $H(z)$):
\[
x[0]=\frac{y[0]}{h[0]},\qquad
x[n]=\frac{1}{h[0]}\Bigl(y[n]-\sum_{k=1}^{n}h[k]\,x[n-k]\Bigr),\quad n\ge 1 .
\]
Also, $\text{length}(x)=\text{length}(y)-\text{length}(h)+1$.

\textbf{(a)} $\{y\}=\{2,5,-5,0,12,10,3\}$ (length 7), $\{h\}=\{2,-3,1,3\}$ (length 4)
$\Rightarrow$ $x$ has length 4:
\[
\begin{aligned}
x[0]&=2/2=1,\\
x[1]&=\tfrac12\bigl(5-(-3)(1)\bigr)=4,\\
x[2]&=\tfrac12\bigl(-5-(-3)(4)-(1)(1)\bigr)=3,\\
x[3]&=\tfrac12\bigl(0-(-3)(3)-(1)(4)-(3)(1)\bigr)=1 .
\end{aligned}
\qquad
\boxed{\{x[n]\}=\{1,\,4,\,3,\,1\}}
\]

\textbf{(b)} $\{y\}=\{0,-15,-7,9,-4,2\}$, $\{h\}=\{5,-1,1\}$ $\Rightarrow$ $x$ has length 4:
\[
\begin{aligned}
x[0]&=0/5=0,\\
x[1]&=\tfrac15\bigl(-15-(-1)(0)\bigr)=-3,\\
x[2]&=\tfrac15\bigl(-7-(-1)(-3)-(1)(0)\bigr)=-2,\\
x[3]&=\tfrac15\bigl(9-(-1)(-2)-(1)(-3)\bigr)=2 .
\end{aligned}
\qquad
\boxed{\{x[n]\}=\{0,\,-3,\,-2,\,2\}}
\]
The two trailing samples of $y$ act as a consistency check:
$y[4]=h[1]x[3]+h[2]x[2]=(-1)(2)+(1)(-2)=-4$ and $y[5]=h[2]x[3]=2$. \checkmark

\textbf{(c)} $\{y\}=\{6,8,-12,-11,14,-22,-12,3,-18\}$ (length 9), $\{h\}=\{-3,2,2,0,3\}$
(length 5) $\Rightarrow$ $x$ has length 5:
\[
\boxed{\{x[n]\}=\{-2,\,-4,\,0,\,1,\,-6\}}
\]
Verification of the last sample: the tail of $h\circledast x$ is $h[4]x[4]=3\cdot(-6)=-18$,
matching $y[8]$.
